% Generated draft: requires mathematical and source review.
\documentclass[11pt]{article}
\usepackage[margin=25mm]{geometry}
\usepackage{fontspec}
\setmainfont{Latin Modern Roman}
\usepackage{amsmath,amssymb,mathtools,unicode-math}
\setmathfont{Latin Modern Math}
\usepackage{enumitem,needspace,xurl,hyperref}
\setlength{\parindent}{0pt}
\setlength{\parskip}{5pt}
\setlength{\emergencystretch}{3em}
\newcommand{\N}{\mathbb{N}}
\newcommand{\Z}{\mathbb{Z}}
\newcommand{\Q}{\mathbb{Q}}
\newcommand{\R}{\mathbb{R}}
\newcommand{\C}{\mathbb{C}}
\newcommand{\F}{\mathbb{F}}
\newcommand{\A}{\mathbb{A}}
\newcommand{\PP}{\mathbb{P}}
\newcommand{\E}{\mathbb{E}}
\newcommand{\eps}{\varepsilon}
\newcommand{\dd}{\,\mathrm{d}}
\begin{document}
\section*{UnsolvedMath ID: 30001023}
\textbf{Source problem code:} catalogue entry
\subsection{Definitions and mathematical statement}
The source compares two continuum frameworks for growth and remodelling of soft biological tissues. In a one-component description, a density $\rho$, velocity $v$ and volumetric mass source $s$ satisfy a balance such as
\[
\partial_t\rho+\nabla\!\cdot(\rho v)=s.
\]
In a mixture description, constituent densities $\rho_\alpha$ and velocities $v_\alpha$ may satisfy separate balances
\[
\partial_t\rho_\alpha+\nabla\!\cdot(\rho_\alpha v_\alpha)
=s_\alpha+q_\alpha,
\]
where internal mass transfers have $\sum_\alpha q_\alpha=0$. Constrained mixtures may impose a common velocity; no particular constituent or constraint class is assumed by the extract. Both frameworks additionally require constitutive mechanics and growth/remodelling laws.

The question is whether mixture theory is a more effective mathematical framework than the usual single-component theory with mass sources for the intended tissue phenomena. A comparison must specify the tissue class, phenomena and observables, model/constitutive families, experimental or benchmark data, and the criterion for effectiveness. Predictive accuracy, identifiability, mathematical well-posedness and computational cost would give different comparisons; none is selected in the source.

Accordingly the record is retained as a modeling-comparison question, not a universal inequality or a theorem that every mixture model outperforms every single-component model. The displayed mass balances explain the distinction without supplying a new constitutive theory or treating mathematical consistency alone as empirical effectiveness.
\subsection{Short English statement}
Is mixture theory more effective than a single-component continuum with volumetric mass sources for soft-tissue growth and remodelling? The source supplies no tissue/model class or effectiveness criterion, so the comparison remains a modeling question.
\subsection{Sources}
\begin{itemize}
\item Supplied report: \emph{Mini-Workshop: The Mathematics of Growth and Remodelling of Soft Biological Tissues} (2009), Oberwolfach Reports, DOI 10.4171/owr/2008/39, pp. 2219--2260; proposed year 2008. \url{https://doi.org/10.4171/owr/2008/39}.
\item Supplied reference: Ambrosi, Ateshian, Humphrey and Oomens, continuum mixture models, \url{https://doi.org/10.1146/annurev-bioeng-071910-124726}.
\item Local UnsolvedMath registry, permanent ID 30001023, catalogue entry; created 25 August 2026, updated 16 September 2026. Its exact-statement verification reports coherence and absence of layout artifacts. Literature checked at \texttt{2026-08-21}: the assessment is supplied dated curation, not independently verified. It reports extensive development/applications of mixtures and labels general comparative effectiveness open as a modeling question rather than a settled theorem.
\end{itemize}
\end{document}
